% paper3_methods.tex -- Materials and Methods of Paper 3, input by
% papers/paper3_intercropping.tex.
\section{Materials and Methods}

\subsection{Treatments and design}
Three cropping systems were compared: sole maize (SOLE\_M), sole cowpea
(SOLE\_C), and an intercrop of the two crops in alternating single rows
(INTER). Sole crops were sown at their recommended densities; in the intercrop
each crop occupied half of the rows. The systems were arranged in randomised
complete blocks with four blocks at each of the three sites
\parencite{gomez1984}, so that every intercrop plot has sole-crop plots of both
species in the same block. Yields are expressed in \si{\tonne\per\hectare} of
dry grain. The data are held in \texttt{data/trials/t3\_intercrop.tsv}; a sole
crop has no yield for the other species, recorded as \texttt{NA}.

\subsection{Land equivalent ratio}
For each block the partial land equivalent ratios of maize and cowpea are
\begin{equation}
  \mathrm{LER}_m = \frac{Y_{im}}{Y_{sm}}, \qquad
  \mathrm{LER}_c = \frac{Y_{ic}}{Y_{sc}},
  \label{eq:partialler}
\end{equation}
where $Y_{im}$ and $Y_{ic}$ are the intercrop yields of maize and cowpea and
$Y_{sm}$ and $Y_{sc}$ their sole-crop yields in the same block. The total land
equivalent ratio is their sum \parencite{mead1980},
\begin{equation}
  \mathrm{LER} = \mathrm{LER}_m + \mathrm{LER}_c .
  \label{eq:ler}
\end{equation}
Standardising each intercrop by the sole crops of its own block keeps the
comparison within the same site, season and soil. An intercrop is said to give
a yield advantage at a site when the mean LER over the four blocks exceeds 1.
